% ===== 第一章 问题重述、符号说明与数据探索（约3页） =====
\section*{一、问题重述与符号说明}

\subsection{问题重述}
“板凳龙”是一种以多节板凳首尾铰接而成的民间舞龙道具。本题中龙头由一节长 341cm 的板凳（含前、后两个把手），龙身与龙尾共 222 节长 220cm 的板凳组成。相邻板凳通过把手销钉铰接，各把手中心沿一条阿基米德螺线分布，螺距记作 $p$。龙头前把手以恒定速率 $1\,\text{m/s}$ 沿螺线向前行进（盘入），其余把手随链运动。总链长（把手中心沿螺线的弧长闭合链）为 $369.16$m。

本题共五问，本文逐一求解：[log:problem]
\begin{enumerate}
\item \textbf{问题一}：螺距 $p=0.55$m，龙头前把手初始位于螺线第 16 圈，求 $t=0\thicksim300$s 内各把手中心的逐秒位置与速度，并列指定节点表；
\item \textbf{问题二}：判定盘入终止（碰撞）时刻；
\item \textbf{问题三}：求使龙头前把手能盘入到以螺线中心为圆心、直径 9m 调头空间的边界（半径 4.5m）的最小螺距；
\item \textbf{问题四}：螺距 1.7m，构造“盘入—S形调头—中心对称盘出”全过程 $-100\thicksim100$s 的位置与速度，并讨论调头圆弧可否缩短；
\item \textbf{问题五}：沿运动路径，求使各把手速度均不超过 2\,m/s 的龙头最大行进速度。
\end{enumerate}

\subsection{符号说明}
本文所用主要符号见表 \ref{tab:sym}。[log:problem, scale]

\begin{table}[H]\centering
\caption{主要符号说明[log:scale]}
\label{tab:sym}\begin{tabular}{lll}
\toprule
符号 & 意义 & 值/单位 \\
\midrule
$p$ & 螺距 & 0.55m(Q1/2)、1.7m(Q4)、0.30m(Q3) \\
$a$ & 螺线尺度 $a=p/2\pi$ & m/rad \\
$r,\theta$ & 螺线极坐标 & m、rad \\
$L_h,L_b$ & 龙头/龙身孔距 & 2.86m、1.65m \\
$w$ & 板凳板宽 & 0.30m \\
$N$ & 把手中心数 & 224 \\
$D_i$ & 把手 $i$ 距龙头前把手的弧长偏移 & 0,…,369.16m \\
$s(\theta)$ & 螺线弧长闭式 & m \\
$S_0$ & 龙头初始弧长 $s(32\pi)$ & 442.59026m \\
$t^{\ast}$ & Q2 终止时刻 & 73.430s \\
$p^{\ast}$ & Q3 最小螺距 & 0.30m \\
$R_1,R_2$ & Q4 S形圆弧半径 & 3.0054m、1.5027m \\
$\beta_{\max}$ & Q5 龙头最大速度 & 2.0\,m/s \\
\bottomrule
\end{tabular}
\end{table}

\subsection{数据探索与自变量确定}
本题无外部数据附件，全部参数由题面几何推导（板长 $341$cm、$220$cm，孔距等于板长减去 $2\times27.5$cm 的孔径偏置，故龙头孔距 $2.86$m、龙身孔距 $1.65$m）。初始龙头前把手位于“第16圈”外端，本文取 $\theta_0=32\pi$（$r_0=8.80$m），该选定在问题一假设中声明。$\theta_0$ 的敏感性在第七章灵敏度分析中给出。独立复核显示，龙头 $r=8.80$m 时，等弧距铺排下龙尾后把手位于 $r\approx3.58$m，初始盘面覆盖螺线第 6.5--16 圈，无自相交。[log:q1.s0]